% ftfp_proposal.tex
\documentclass[hyperref, xcolor=dvipsnames]{beamer}

\usepackage{amsmath, caption, subcaption, graphicx, tikz}
\title[Thesis Proposal Defense]{Approximation Algorithms for the
  Facility Location prolems}
\author[lyan]{Li Yan}
\institute[UCR]{
  Computer Science\\
  University of California Riverside\\
  [\medskipamount]
  \includegraphics[width=.2\textwidth,height=.25\textheight]{figure/UCseal.eps}
}
\date{Jan 30th, 2013}
\usetheme{Warsaw}


%%%%%%%%%%%%%%%%%%%%%%%%%%%%%%%%%%%%%%%%%%%%%%%%%%%%%%%%%%

% non-math stuff

\newcommand{\myparagraph}[1]{{\smallskip\noindent{\bf #1}}}
\newcommand{\emparagraph}[1]{{\smallskip\noindent{\it #1}}}
\newcommand{\etal}{{\it et al.}}
\newcommand{\myif}{{\mbox{\rm\ if \ }}}
\newcommand{\mycase}[1]{\mbox{{\underline{Case #1}}:\/}}

\newcommand{\margincomment}[1]%
    {{%
      \marginpar{{\tiny\begin{minipage}{0.5in}
                       \begin{flushleft}
                          {#1}
                       \end{flushleft}
                       \end{minipage}
                }}
    }}


%%%%%%%%%%%%%%%%%%%%%%%%%%%%%%%%%%%%%%%%%%%%%%%%%%%%%%%%%%

% various letters

\newcommand{\hatc}{{\hat c}}
\newcommand{\hatC}{{\hat C}}
\newcommand{\hatr}{{\hat r}}
\newcommand{\hatx}{{\hat x}}
\newcommand{\haty}{{\hat y}}
\newcommand{\dotx}{{\dot x}}
\newcommand{\doty}{{\dot y}}
\newcommand{\dotr}{{\dot r}}
\newcommand{\boldx}{{\mathbf x}}

\newcommand{\doubledone}{{\bar 1}}
\newcommand{\doubledtwo}{{\bar 2}}
\newcommand{\barc}{{\bar c}}
\newcommand{\bart}{{\bar t}}

\newcommand{\barx}{{\bar x}}
\newcommand{\bary}{{\bar y}}
\newcommand{\barz}{{\bar z}}
\newcommand{\barr}{{\bar r}}
\newcommand{\barX}{{\bar X}}
\newcommand{\barY}{{\bar Y}}
\newcommand{\barZ}{{\bar Z}}
\newcommand{\bara}{{\bar a}}
\newcommand{\bard}{{\bar d}}
\newcommand{\barm}{{\bar m}}
\newcommand{\barA}{{\bar A}}
\newcommand{\barB}{{\bar B}}
\newcommand{\barC}{{\bar C}}
\newcommand{\barG}{{\bar G}}
\newcommand{\barE}{{\bar E}}
\newcommand{\barV}{{\bar V}}

\newcommand{\wbarC}{{\overline{C}}}
\newcommand{\wbarD}{{\overline{D}}}
\newcommand{\wbarN}{{\overline{N}}}
\newcommand{\wbarX}{{\overline{X}}}


\newcommand{\barbeta}{{\bar\beta}}
\newcommand{\bargamma}{{\bar\gamma}}
\newcommand{\apomega}{{\bar\omega}}

\newcommand{\bfr}{\boldsymbol{r}}
\newcommand{\bfv}{{\bf v}}
\newcommand{\bfx}{\boldsymbol{x}}
\newcommand{\bfy}{\boldsymbol{y}}
\newcommand{\bfz}{{\bf z}}
\newcommand{\bfQ}{{\bf Q}}
\newcommand{\bfR}{{\bf R}}
\newcommand{\bfS}{{\bf S}}
\newcommand{\bfT}{{\bf T}}
\newcommand{\bfV}{{\bf V}}
\newcommand{\bfone}{{\bf 1}}
\newcommand{\bfalpha}{\boldsymbol{\alpha}}
\newcommand{\bfbeta}{\boldsymbol{\beta}}

\newcommand{\calA}{{\cal A}}
\newcommand{\calB}{{\cal B}}
\newcommand{\calC}{{\cal C}}
\newcommand{\calD}{{\cal D}}
\newcommand{\calE}{{\cal E}}
\newcommand{\calG}{{\cal G}}
\newcommand{\calH}{{\cal H}}
\newcommand{\calJ}{{\cal J}}
\newcommand{\calK}{{\cal K}}
\newcommand{\calL}{{\cal L}}
\newcommand{\calM}{{\cal M}}
\newcommand{\calN}{{\cal N}}
\newcommand{\calS}{{\cal S}}
\newcommand{\calU}{{\cal U}}
\newcommand{\calX}{{\cal X}}
\newcommand{\calT}{{\cal T}}

\newcommand{\hatcalI}{{\hat{\cal I}}}
\newcommand{\barcalI}{{\bar{\cal I}}}
\newcommand{\dotcalI}{{\dot{\cal I}}}

\newcommand{\vecS}{{\bar S}}
\newcommand{\vecT}{{\bar T}}
\newcommand{\vecone}{{\bf 1}}
\newcommand{\tildec}{{\tilde c}}
\newcommand{\tilded}{{\tilde d}}
\newcommand{\tildeD}{{\tilde D}}
\newcommand{\tildeC}{{\widetilde C}}
\newcommand{\tildeZ}{{\tilde Z}}
\newcommand{\tilder}{{\widetilde r}}
\newcommand{\tildex}{{\widetilde x}}
\newcommand{\wtildeN}{{\widetilde N}}
\newcommand{\tildebfr}{\widetilde{\boldsymbol{r}}}
\newcommand{\tildebfx}{\widetilde{\boldsymbol{x}}}
\newcommand{\tildebfy}{\widetilde{\boldsymbol{y}}}

\newcommand{\barbfx}{\bar{\boldsymbol{x}}}
\newcommand{\barbfy}{\bar{\boldsymbol{y}}}
\newcommand{\hatbfx}{\hat{\boldsymbol{x}}}
\newcommand{\hatbfy}{\hat{\boldsymbol{y}}}
\newcommand{\dotbfx}{\dot{\boldsymbol{x}}}
\newcommand{\dotbfy}{\dot{\boldsymbol{y}}}

\newcommand{\wbarcalC}{{\overline{\calC}}}
\newcommand{\wbarcalD}{{\overline{\calD}}}
\newcommand{\eps}{{\epsilon}}

%%%%%%%%%%%%%%%%%%%%%%%%%%%%%%%%%%%%%%%%%%%%%%%%%%%%%%%%%%

\newcommand{\half}{{\mbox{$\frac{1}{2}$}}}
\newcommand{\threehalfs}{{\mbox{$\frac{3}{2}$}}}
\newcommand{\threefourths}{{\mbox{$\frac{3}{4}$}}}
\newcommand{\fivehalfs}{{\mbox{$\frac{5}{2}$}}}
\newcommand{\onethird}{{\mbox{$\frac{1}{3}$}}}
\newcommand{\twothirds}{{\mbox{$\frac{2}{3}$}}}
\newcommand{\fourthirds}{{\mbox{$\frac{4}{3}$}}}
\newcommand{\fivethirds}{{\mbox{$\frac{5}{3}$}}}
\newcommand{\fivefourths}{{\mbox{$\frac{5}{4}$}}}
\newcommand{\onefourth}{{\mbox{$\frac{1}{4}$}}}
\newcommand{\onefifth}{{\mbox{$\frac{1}{5}$}}}
\newcommand{\twofifths}{{\mbox{$\frac{2}{5}$}}}
\newcommand{\threefifths}{{\mbox{$\frac{3}{5}$}}}
\newcommand{\fourfifths}{{\mbox{$\frac{4}{5}$}}}
\newcommand{\ninefifths}{{\mbox{$\frac{9}{5}$}}}
\newcommand{\sevensixths}{{\mbox{$\frac{7}{6}$}}}
\newcommand{\oneeighth}{{\mbox{$\frac{1}{8}$}}}
\newcommand{\threeeighths}{{\mbox{$\frac{3}{8}$}}}
\newcommand{\fiveeighths}{{\mbox{$\frac{5}{8}$}}}
\newcommand{\seveneighths}{{\mbox{$\frac{7}{8}$}}}
\newcommand{\onetenth}{{\mbox{$\frac{1}{10}$}}}
\newcommand{\seventenths}{{\mbox{$\frac{7}{10}$}}}
\newcommand{\ninetenths}{{\mbox{$\frac{9}{10}$}}}
\newcommand{\twonineths}{{\mbox{$\frac{2}{9}$}}}
\newcommand{\fivenineths}{{\mbox{$\frac{5}{9}$}}}
\newcommand{\elevennineths}{{\mbox{$\frac{11}{9}$}}}
\newcommand{\threetwentieths}{{\mbox{$\frac{3}{20}$}}}
\newcommand{\twentyfivenineteenths}{{\mbox{$\frac{25}{19}$}}}

\newcommand{\sqrttwo}{\sqrt{2}}

%%%%%%%%%%%%%%%%%%%%%%%%%%%%%%%%%%%%%%%%%%%%%%%%%%%%%%%%%%

% various delimiters

\newcommand{\braced}[1]{{ \left\{ #1 \right\} }}
\newcommand{\angled}[1]{{ \left\langle #1 \right\rangle }}
\newcommand{\brackd}[1]{{ \left[ #1 \right] }}
\newcommand{\parend}[1]{{ \left( #1 \right) }}
\newcommand{\barred}[1]{{ \left| #1 \right| }}
\newcommand{\dbarred}[1]{{ \left\| #1 \right\| }}
\newcommand{\floor}[1]{{ \lfloor #1 \rfloor }}
\newcommand{\ceiling}[1]{{ \lceil #1 \rceil }}

%%%%%%%%%%%%%%%%%%%%%%%%%%%%%%%%%%%%%%%%%%%%%%%%%%%%%%%%%%

% some math symbols

\newcommand{\set}{\,{\leftarrow}\,}
\newcommand{\suchthat}{{\,:\,}}
\newcommand{\cost}{{\it cost}}
\newcommand{\yield}{{\it yield}}
\newcommand{\opt}{{\it opt}}

\newcommand{\algA}{{\bf A}}
\newcommand{\LHS}{{\rm LHS}}
\newcommand{\RHS}{{\rm RHS}}
\newcommand{\reals}{{\bf R}}
\newcommand{\posreals}{{\bf R}^+}

\newcommand{\assign}{{\,\leftarrow\,}}

\newcommand{\absvalue}[1]{{\barred{#1}}}
\newcommand{\posvalue}[1]{{\brackd{#1}^+}}

\newcommand{\NP}{{\mbox{\sf NP}}}
\newcommand{\PP}{{\mbox{\sf P}}}
\newcommand{\DTIME}{{\mbox{\sf DTIME}}}

\newcommand{\letbox}[1]{{\makebox[11pt]{{\small {$#1$}}}}}
\newcommand{\optstring}[1]{{ \frame{\;\raisebox{0pt}[12pt][5pt]{#1}\;} }}

\newcommand{\leftend}{{\diamond}}
\newcommand{\rightend}{{\diamond}}

%\newcommand{\argmin}{{\mbox{\rm argmin}}}
\DeclareMathOperator*{\argmin}{arg\,min}

\newcommand\litem[1]{\item{\bfseries #1\enspace}}
\newcommand{\ceil}[1] {\lceil #1 \rceil}
\newcommand{\naive}{na\"{\i}ve}
\newcommand{\LP}{\mbox{\rm LP}}
\newcommand{\OPT}{\mbox{\rm OPT}}
\newcommand{\ALG}{\mbox{\rm ALG}}
\newcommand{\LPR}[1]{{\mbox{\rm LPR#1}}}
\newcommand{\smallLPR}[1]{{\mbox{\tiny\rm LPR#1}}}
% algorithm names
\newcommand{\ESTA}{\mbox{\rm ESTA}} % 4approx
\newcommand{\EGUP}{\mbox{\rm EGUP}} % 3approx
\newcommand{\ECHS}{\mbox{\rm ECHS}} % 1.736
\newcommand{\EBGS}{\mbox{\rm EBGS}} % 1.575
\newcommand{\GUP}{\mbox{\rm GUP}}
\newcommand{\smallESTA}{\mbox{\tiny\rm ESTA}}
\newcommand{\smallEGUP}{\mbox{\tiny\rm EGUP}}
\newcommand{\smallECHS}{\mbox{\tiny\rm ECHS}}
\newcommand{\smallEBGS}{\mbox{\tiny\rm EBGS}}

\newcommand{\SOL}[1]{{{\mbox{\rm SOL}}_{#1}}}
\newcommand{\FTFP}{\mbox{\rm FTFP}}
\newcommand{\FTFL}{\mbox{\rm FTFL}}
\newcommand{\calI}{\mathcal{I}}
\newcommand{\avg}{{\mbox{\scriptsize\rm avg}}}

\newcommand{\dmax}{\text{dmax}}
\newcommand{\davg}{\text{davg}}
\newcommand{\favg}{f_{\text{avg}}}
\newcommand{\conn}{\text{conn}}
\newcommand{\cls}{\text{cls}}
\newcommand{\far}{\text{far}}

\newcommand{\sitesset}{\mathbb{F}}
\newcommand{\clientset}{\mathbb{C}}
\newcommand{\facilityset}{\overline{\sitesset}}
\newcommand{\demandset}{\overline{\clientset}}

%\newcommand{\dist}{{\mbox{dist}}}
\newcommand{\concost}{C^{\avg}}
\newcommand{\faccost}{F^{\avg}}
\newcommand{\tcc}{\textrm{tcc}}
\newcommand{\clsdist}{C_{\cls}^{\avg}}
\newcommand{\fardist}{C_{\far}^{\avg}}
\newcommand{\clsmax}{C_{\cls}^{\max}}
\newcommand{\clsnb}{N_{\cls}}
\newcommand{\farnb}{N_{\far}}
\newcommand{\wbarclsnb}{\wbarN_{\cls}}
\newcommand{\wbarfarnb}{\wbarN_{\far}}

\newcommand{\Exp}{\mbox{\rm Exp}}

\newcommand{\FacilityDistSort}{{\textsc{FacilityDistSort}}}
\newcommand{\NearestUnitChunk}{{\textsc{NearestUnitChunk}}}
\newcommand{\AugmentToUnit}{{\textsc{AugmentToUnit}}}
\newcommand{\connsum}{{\textrm{conn}}}

%%%%%%%%%%%%%%%%%%%%%%%%%%%%%%%%%%%%%%%%%%%%%%%%%%%%%%%%%%

% theorem and such

\newtheorem{fact}[theorem]{Fact}
\newtheorem{observation}[theorem]{Observation}

%%%%%%%%%%%%%%%%%%%%%%%%%%%%%%%%%%%%%%%%%%%%%%%%%%%%%%%%%%

\newcommand{\ignore}[1]{}

\begin{document}

%%%
\begin{frame}
  \titlepage
\end{frame}

%%%
\begin{frame}
  \frametitle{The Uncapacitated Facility Location Problem (UFL)}
  All demands are 1, each facility can be opened only once.
  \begin{tikzpicture}[scale=0.5]
    \node[fill=blue!50] (fac1) at (0,2) {$f_1$};
    \node[fill=blue!50] (fac2) at (0,5) {$f_2$};
    \node[fill=blue!50] (fac3) at (0,8) {$f_3$};

    \node[circle,fill=green!50] (client1) at (6,0) {$r_1=1$};
    \node[circle,fill=green!50] (client2) at (6,3) {$r_2=1$};
    \node[circle,fill=green!50] (client3) at (6,6) {$r_3=1$};
    \node[circle,fill=green!50] (client4) at (6,9) {$r_4=1$};
    
    \foreach \from in {fac1,fac2,fac3}
    \foreach \to in {client1,client2,client3,client4}
        \draw (\from)--(\to);
  \end{tikzpicture}

\end{frame}

%%%
\begin{frame}
  \frametitle{The Fault-tolerant Facility Location Problem (FTFL)}
  Demands may be more than 1, each facility can be opened only once.

  \begin{tikzpicture}[scale=0.5]
    \node[fill=blue!50] (fac1) at (0,2) {$f_1$};
    \node[fill=blue!50] (fac2) at (0,5) {$f_2$};
    \node[fill=blue!50] (fac3) at (0,8) {$f_3$};

    \node[circle,fill=green!50] (client1) at (6,0) {$r_1$};
    \node[circle,fill=green!50] (client2) at (6,3) {$r_2$};
    \node[circle,fill=green!50] (client3) at (6,6) {$r_3$};
    \node[circle,fill=green!50] (client4) at (6,9) {$r_4$};
    
    \foreach \from in {fac1,fac2,fac3}
    \foreach \to in {client1,client2,client3,client4}
        \draw (\from)--(\to);
  \end{tikzpicture}
\end{frame}

%%%
\begin{frame}
  \frametitle{The Fault-tolerant Facility Placement Problem (FTFP)}
  Demands may be more than 1, each facility can be opened multiple times.

  \begin{tikzpicture}[scale=0.5]
    \node[fill=blue!50] (fac1) at (0,2) {$f_1$};
    \node[fill=blue!50] (fac2) at (0,5) {$f_2$};
    \node[fill=blue!50] (fac3) at (0,8) {$f_3$};

    \node[circle,fill=green!50] (client1) at (6,0) {$r_1$};
    \node[circle,fill=green!50] (client2) at (6,3) {$r_2$};
    \node[circle,fill=green!50] (client3) at (6,6) {$r_3$};
    \node[circle,fill=green!50] (client4) at (6,9) {$r_4$};
    
    \foreach \from in {fac1,fac2,fac3}
    \foreach \to in {client1,client2,client3,client4}
        \draw (\from)--(\to);
  \end{tikzpicture}

\end{frame}

%%%
\begin{frame}
  \frametitle{Best Known Approximation Results}
  \begin{itemize}
  \item UFL: 1.488, a combination of LP-rounding and greedy, by Li (Princeton)
  \item FTFL: 1.7245, dependent rounding and laminar clustering, by Byrka, Srinivasan and Swamy (U Maryland)
  \item FTFP: 1.575, LP-rounding (UCR)
  \end{itemize}
\end{frame}

%%%
\begin{frame}
  \frametitle{Lower Bound on Approximability}
  \begin{itemize}
  \item No ratio better than $1.463$ unless $\PP = \NP$. Reduction
    from Set Cover, by Sviridenko.
  \item Integrality Gap is also $1.463$, the example uses $n$ facilities and $n \choose l$ clients. The fractional solution is each $y_i = 1/l$.
  \end{itemize}
  \begin{tikzpicture}[scale=.7,rotate=-90]
    \node[fill=blue!50] at (.5,2) {$f_1$};
    \node[fill=blue!50] at (.5,5) {$f_2$};
    \node[fill=blue!50] at (.5,13) {$f_n$};
    
    \fill (.5,7.5) circle (5pt);
    \fill (.5,9) circle (5pt);
    \fill (.5,10.5) circle (5pt);

    \node[above] (fac1) at (0,2) {$y_1=1/l$};
    \node[above] (fac2) at (0,5) {$y_2=1/l$};
    \node[above] (fac3) at (0,13) {$y_n=1/l$};

    \draw[fill=green!50] (2.5,0) circle (.5);
    \draw[fill=green!50] (2.5,3) circle (.5);
    \draw[fill=green!50] (2.5,6) circle (.5);
    \draw[fill=green!50] (2.5,15) circle (.5);

    \fill (2.5,8) circle (5pt);
    \fill (2.5,10) circle (5pt);
    \fill (2.5,12) circle (5pt);

    \node[below] (client1) at (3,0) {$1$};
    \node[below] (client2) at (3,3) {$2$};
    \node[below] (client3) at (3,6) {$3$};
    \node[below] (client4) at (3,15) {$n \choose l$};
  \end{tikzpicture}
\end{frame}

%%%
\begin{frame}
  \frametitle{The LP Formulation}
  \begin{itemize}
  \item $y_i$ represent the number of facilities built at site $i$.
  \item $x_{ij}$ represent the number of connections from client $j$
    to facilities at site $i$.
  \end{itemize}
  %%%%%%%%%%% 
  \begin{alignat}{3}
    \textrm{minimize}\quad \textstyle{\sum_{i\in \sitesset}}f_iy_i &+ \textstyle{\sum_{i\in \sitesset, j\in \clientset}}d_{ij}x_{ij}&\label{eqn:fac_primal}
    \\ \notag
    \textrm{subject to}\quad y_i - x_{ij} &\geq 0  & &\forall i\in \sitesset, j\in \clientset 
    \\ \notag
    \textstyle{\sum_{i\in \sitesset} x_{ij}} &\geq r_j & &\forall j\in \clientset
    \\ \notag
    x_{ij} \geq 0, y_i &\geq 0 & &\forall i\in \sitesset, j\in \clientset 
    \\ \notag
  \end{alignat}
  %%%%%%%%%%%%
  \begin{alignat}{3}
    \textrm{maximize}\quad \textstyle{\sum_{j\in \clientset}} r_j\alpha_j&\label{eqn:fac_dual}  
    \\ \notag
    \textrm{subject to} \quad \textstyle{
      \sum_{j\in \clientset}\beta_{ij}} &\leq f_i  &\quad\quad			&\forall i \in \sitesset  
    \\ \notag
    \alpha_{j} - \beta_{ij} 	&\leq  d_{ij}       &                 & \forall i\in \sitesset, j\in \clientset 
    \\ \notag
    \alpha_j \geq 0, \beta_{ij} &\geq 0           &            & \forall i\in \sitesset, j\in \clientset
  \end{alignat}
  %%%%%%%%%%%
\end{frame}

%%%
\begin{frame}
  \frametitle{The Shmoys, Tardos and Ardal's Algorithm (STA97)}
  \begin{itemize}
  \item If all $N(j)$ disjoint, then easy.
  \item To bound $F^A$, need neighborhood of chosen clients be disjoint.
  \item To bound $C^A$, need non-primary clients having a fail over connection.
  \end{itemize}

  \begin{itemize}
  \item The greedy clustering: iteratively find the best client and assign other clients to it.
  \item Estimate $\max d_{ij}$, either cut the neighborhood $N(j)$ or use dual solution.
  \end{itemize}
\end{frame}

%%%
\begin{frame}
  \frametitle{Chudak, Svi, Byrka and Li' improvement}
  \begin{itemize}
  \item Chudak: randomized rounding, estimate on the expected connection cost
  \item Sviridenko: use a concave function to upper bound distance and guide rounding
  \item Byrka: boost facility opening probability and use $N_{\cls}(j)$ for overlapping
  \item Li: find the right distribution for probability boost
  \end{itemize}
\end{frame}

%%%
\begin{frame}
  \frametitle{Demand Reduction} Given an FTFP instance, we can reduce
  it to an instance such that $R=\max_{j}r_j$ is bounded by
  $2|\sitesset|$.

  \begin{itemize}
    \item $\hatx_{ij} = \floor{x_{ij}^\ast}, \haty_{i} = \floor{y_i^\ast}$
    \item $\dotx_{ij} = x_{ij}^\ast - \hatx_{ij}, \doty_{i} = y_i^\ast
      - \haty_{i}$
    \item $\hatr_j = \sum_{i\in\sitesset} \hatx_{ij}$
    \item $\dotr_j = r_j - \hatr_j$
  \end{itemize}
  
  \begin{claim}
    $\hatx_{ij}, \haty_{i}$ is feasible and optimal for $\hat I$, and
    $\dotx_{ij}, \doty_{i}$ is feasible and optimal for $\dot I$.
  \end{claim}
\end{frame}

%%%
\begin{frame}
  \begin{theorem}
    Given any $\rho\geq 1$ approximation algorithm $\calA$ for solving
    restricted FTFP, we can obtain an algorithm with
    $\rho$-approximation for general FTFP.
  \end{theorem}
  \begin{proof}
    Solve LP and obtain $\hat \calI$ and $\dot \calI$. For $\hat I$ we have
    ratio $1$, and use $\calA$ to solve $\dot \calI$ with ratio
    $\rho$. Final ratio is $\max\{1,\rho\}$.
  \end{proof}
\end{frame}

%%%
\begin{frame}
  \begin{corollary}
    There is a $1.7245$-approximation algorithm for the FTFP problem.
  \end{corollary}
  \begin{proof}
    We simply reduce an FTFP problem into an FTFL problem. Solving the
    instance $\dot \calI$ using the $1.7245$-approximation algorithm
    for FTFL by Byrka {\etal}.
  \end{proof}
\end{frame}

%%%
\begin{frame}
  \frametitle{Improve from $1.7245$ to $1.575$-approximation}
  \begin{itemize}
  \item We have shown that FTFP can be reduced to FTFL while
    preserving the approximation ratio.
  \item Next step is to show FTFP can be approximated with a better
    ratio than FTFL.
  \item Simple case is when all $r_j$'s are equal, then we can apply
    any UFL approximation results to FTFP as the uniform FTFP is
    simply a scaled version of UFL.
  \item For general FTFP, we need \emph{Adaptive Partition}.
  \end{itemize}
\end{frame}

%%%
\begin{frame}
  \frametitle{Adaptive Partition} 

  Given an instance of FTFP, with its fractional optimal solution
  $(\bfx^\ast, \bfy^\ast)$, w.l.o.g. we assume \emph{completeness},
  i.e. $x_{ij}^\ast > 0$ implies $x_{ij} ^\ast = y_i^\ast$.  Then we
  can partition the instance into unit demands and facilities, with
  fractional solution $(\barbfx, \barbfy)$ such that
  \begin{itemize}
  \item $x_{ij}^\ast$ is spread among its demands.
  \item $y_i^\ast$ is spread among its facilities.
  \item Each demand $\nu$ has a neighborhood $\wbarN(\nu)$ with total
    connectoin value of $1$.
  \item Primary demands have a smaller cost than non-primary demands
    assigned to them.
  \end{itemize}
\end{frame}

%%%
\begin{frame}
  \frametitle{An Example of Adaptive Partition}
  4 sites and 4 clients,
  \begin{itemize}
  \item all $f_i=1$
  \item all $d_{ij}=1$ except $d_{ii} = 3$ for $i=1,2,3,4$
  \item $r_1 = 1, r_2 = r_3 = r_4 = 2$.
  \end{itemize}
\end{frame}

%%%
\begin{frame}
  \frametitle{The Fractional Optimal Solution}
  \begin{table}
    \begin{subtable}{0.5\textwidth}
      \centering
      \begin{tabular}{c | c c c c}
        $y_i^\ast$ & 1 & 2 & 3 & 4\\
        \hline
        & 4/3 & 1/3 & 1/3 & 1/3\\
      \end{tabular}
      \subcaption{}
    \end{subtable}
%
    \begin{subtable}{0.5\textwidth}
      \centering
      \begin{tabular}{c | c c c c}
        $x_{ij}^\ast$ & $j=1$ & 2 & 3 & 4\\
        \hline
        $i=1$ & 0 & 4/3 & 4/3 & 4/3\\
        2 & 1/3 & 0 & 1/3 & 1/3\\
        3 & 1/3 & 1/3 & 0 & 1/3\\
        4 & 1/3 & 1/3 & 1/3 & 0\\
      \end{tabular}
      \subcaption{}
    \end{subtable}
    \caption{An optimal fractional solution to the FTFP instance.}
\end{table}

The dual solution has all $\alpha_j^\ast = 4/3$.
\end{frame}

%%%
\begin{frame}
  \frametitle{Phase 1: Iteration 1} 

  Choose client 1 and create a primary demand $\kappa_{11}$. Each of
  client 2,3,4 creates a demand and assigned to $\kappa_{11}$.
  \begin{table}
  \begin{subtable}{0.5\textwidth}
    \centering
    \begin{tabular}{c | c c c c}
      $\bary$ & 1 & 2 & 3 & 4\\
      \hline
      & 1/3 & 1/3 & 1/3 & 1/3\\
    \end{tabular}
    \subcaption{Opening fraction of the first facility in each site.}
  \end{subtable}
%
  \begin{subtable}{0.5\textwidth}
    \centering
    \begin{tabular}{c | c c c c}
      $\barx$ & 1 & 2 & 3 & 4\\
      \hline
      1 & 0 & 1/3 & 1/3 & 1/3\\
      2 & 0 & 0 & 1/3 & 1/3\\
      3 & 0 & 1/3 & 0 & 1/3\\
      4 & 0 & 1/3 & 1/3 & 0\\
    \end{tabular}
    \subcaption{Connection fraction of the first demand of each client.}
  \end{subtable}
\end{table}
\end{frame}

\begin{frame}
  \frametitle{Phase 1: Iteration 2}

  Choose client 4 and create a primary demand $\kappa_{42}$. Each of
  client 2,3 creates a demand and assigned to $\kappa_{42}$.
  \begin{table}
  \begin{subtable}{0.5\textwidth}
    \centering
    \begin{tabular}{c | c c c c}
      $\bary$ & 1 & 2 & 3 & 4\\
      \hline
      & 1 &  &  & \\
    \end{tabular}
    \subcaption{Opening fraction of the second facility in each site.}
  \end{subtable}
%
  \begin{subtable}{0.5\textwidth}
    \centering
    \begin{tabular}{c | c c c c}
      $\barx$ & 1 & 2 & 3 & 4\\
      \hline
      2 & 1 & 0 & 0 & 0\\
      3 & 1 & 0 & 0 & 0\\
      4 & 1 & 0 & 0 & 0\\
    \end{tabular}
    \subcaption{Connection fraction of the second demand of each
      client. Notice client 1 does not have a second demand.}
  \end{subtable}
\end{table}
\end{frame}

%%%
\begin{frame}
  \frametitle{Phase 2: Augment to Unit}
  \begin{table}
    \begin{subtable}{0.5\textwidth}
      \centering
      \begin{tabular}{c | c c c c}
        $\barx$ & 1 & 2 & 3 & 4\\
        \hline
        1 & 0 & 1/3 & 1/3 & 1/3\\
        2 & 1/3 & 0 & 1/3 & 1/3\\
        3 & 1/3 & 1/3 & 0 & 1/3\\
        4 & 1/3 & 1/3 & 1/3 & 0\\
      \end{tabular}
      \subcaption{Connection fraction of the first demand of each
        client.}
    \end{subtable}
%
    \begin{subtable}{0.5\textwidth}
      \centering
      \begin{tabular}{c | c c c}
        $\barx$ & 2 & 3 & 4\\
        \hline
        1  & 1 & 1 & 1\\
        2  & 0 & 0 & 0\\
        3  & 0 & 0 & 0\\
        4  & 0 & 0 & 0\\
      \end{tabular}
      \subcaption{Connection fraction of the secon demand of each
        client except client 1.}
    \end{subtable}
    \caption{Notice all demands have connection value $1$.}
  \end{table}
\end{frame}

\begin{frame}
  \frametitle{$3$-approximation Algorithm}
  \begin{itemize}
  \item Each primary demand open $\mu\in \wbarN(\kappa)$ with probability
    $\bary_{\mu}$.
  \item Each primary demand connects to the only open facility
    $\phi(\kappa)$ in $\wbarN(\kappa)$.
  \item Each non-primary demand connects to $\phi(\kappa)$.
  \item Expected facility cost at most $F^\ast$.
  \item Expected connection cost at most $C^\ast + 2\;\LP^\ast$.
  \end{itemize}
\end{frame}

\begin{frame}
  \frametitle{$1.736$-approximation Algorithm}
  \begin{itemize}
  \item Improve connection cost estimate: For non-primary demands use
    $\mu$ in $\wbarN(\nu)$ if one is open.
  \item Expected facility cost at most $F^\ast$.
  \item Expected connection cost at most $C^\ast + 2/e\;\LP^\ast$.
  \end{itemize}
\end{frame}

\begin{frame}
  \frametitle{$1.575$-approximation}
  \begin{itemize}
  \item Need a more refined partition to deal with close and far
    neighborhood.
  \item $\wbarN_{\cls}(\nu)$ has total connection value $1/\gamma$.
  \item $\wbarN_{\far}(\nu)$ has total connection value $1-1/\gamma$.
  \item Assignment implies overlap of $\wbarN_{\cls}$ of $\nu$ and
    $\kappa$.
  \item Expected facility cost at most $\gamma F^\ast$.
  \item Expected connection cost at most
    $\max\{\frac{1+1/e^\gamma}{1-1/e^\gamma}, 1+2/e^\gamma\}\; C^\ast$.
  \item Picking $\gamma=1.575$ gives ratio $1.575$.
  \end{itemize}
\end{frame}

\begin{frame}
  \frametitle{Primal-dual Algorithms}
  \begin{itemize}
  \item A Simple $O((\log R/\log\log R)^2)$ Algorithm.
  \item Greedy Algorithm with Dual-fitting Analysis.
  \end{itemize}
\end{frame}

\begin{frame}
  \frametitle{A Simple $O((\log R/\log\log R)^2)$ Algorithm}
  \begin{itemize}
  \item Let $r_1,\ldots, r_n$ be demands of the $n$ clients.
  \item Group clients by $[k^i, k^{i+1})$ for $k$ such that $k^k = R =
    \max_j r_j$.
  \item Solve each group by treating each client with
    $r_j=k^{i+1}$.
  \item Combine all solution to each group to obtain final integral
    solution.
  \end{itemize}
\end{frame}

\begin{frame}
  \frametitle{Performance Analysis}
  \begin{theorem}
    There is a primal-dual $O((\log R /\log\log R)^2)$-approximation
    algorithm for FTFP.
  \end{theorem}
  \begin{proof}
    \begin{itemize}
    \item When solving each group individually, we pay a factor of $k$
      for each group, since within each group the total cost is up by
      at most $k$ times due to $r_j$ rounded up.
    \item When combining solutions, we pay a factor of $k$ since each
      facility can be over counted at most $k$ times, since we have at
      most $k$ groups because $k^k = R$.
    \end{itemize}
  \end{proof}
\end{frame}

\begin{frame}
  \frametitle{The Greedy Algorithm}
  \begin{itemize}
  \item Repeatedly picking the star with minimum average cost.
  \item A star is a facility $i$ and a set of clients $S$.
  \item Average cost is $ (f_i + \sum_{j\in S} d_{ij}) / |S|$.
  \end{itemize}
\end{frame}

\begin{frame}
  \frametitle{Known Facts about Greedy}
  \begin{itemize}
  \item Runs in polynomial time as the best star remains best until
    exhausted so can combine iterations into phases.
  \item $O(H_n)$-approximation by dual-fitting analysis.
  \item Open question: Is it $O(1)$-approximation?
  \end{itemize}
\end{frame}

\begin{frame}
  \frametitle{Summary} We studied the fault-tolerante facility
  placement problem (FTFP) on approximation algorithms.

  \begin{itemize}
  \item Known results
    \begin{itemize}
    \item LP-rounding algorithms achieve a best ratio of $1.575$,
      matching the best LP-based ratio for its special case, UFL.
    \item Primal-dual algorithms achieve $O(\log R/\log\log R)$,
      better than the $O(\log R)$ ratio for primal-dual algorithm for
      FTFL.
    \item The greedy algorithm has ratio no more than $O(H_n)$.
  \end{itemize}
  \item Work in progress:
    Resolve whether Greedy is $O(1)$-approximation or not.
  \end{itemize}
\end{frame}
%%%%%%%%%%
\end{document}
